% ===== 核心排版设置 =====
\usepackage[top=2.5cm, bottom=2.5cm, left=3cm, right=2.5cm, 
            headheight=15pt, includehead]{geometry}
\usepackage{setspace}
\usepackage{titlesec}

% ===== 视觉增强宏包 =====
\usepackage{fontspec}
\usepackage[perpage,symbol*]{footmisc}
\usepackage[dvipsnames, svgnames, x11names]{xcolor}

% ===== 内容结构宏包 =====
\usepackage{array} % for better column formatting
\usepackage{booktabs} % for nicer table rules
\usepackage[most]{tcolorbox}
\usepackage{enumitem}       % 列表定制
\usepackage{graphicx}       % 插图支持
\usepackage{amssymb}         % 示例文本

% ===== 字体配置 =====
\setmainfont{Times New Roman}
\setsansfont{Arial}
\setCJKmainfont[BoldFont=STZhongsong]{SimSun}
\setCJKsansfont{Microsoft YaHei}
\setCJKmonofont{NSimSun}

\definecolor{lessonblue}{RGB}{45, 85, 155}
\definecolor{examplebg}{RGB}{245, 248, 252}

% ===== 例句环境 =====
\tcbset{
  enhanced,
  breakable,
  sharp corners,
  boxsep=6pt,
  colback=examplebg,
  colframe=examplebg,
}

\newtcolorbox{liju}{
  breakable,
  colback=examplebg,
  colframe=examplebg,
%   borderline horizontal={1.2pt}{0pt}{lessonblue}, % horizontal rule
  before upper={%
    % \vspace{0.3\baselineskip}%
    \noindent\textbf{\color{lessonblue}例句}%
    % \vspace{0.3\baselineskip}%
    \par\noindent{\color{lessonblue}\rule{\linewidth}{1.2pt}}%
    \par\vspace{0.3\baselineskip}%
    \setlength{\parindent}{0pt}%
    \setlist[enumerate]{noitemsep, topsep=0pt, leftmargin=*}%
  },
%   after upper={\vspace{0.4\baselineskip}}
}

% 例句内容命令
\newcommand{\enzh}[2]{%
  \noindent \makebox[1em][l]{\raisebox{0.25ex}{\textcolor{lessonblue}{\large\ensuremath{\bullet}}}}%
  {#1}\par
  \noindent\hspace*{1.5em}{\small\textcolor{gray!30!black}{#2}}\par
  \vspace{0.3\baselineskip}%
}

% 正误对照命令
\newcommand{\compare}[2]{%
  % 错误句子行
  \noindent\makebox[1.8em][l]{\textcolor{red!70!black}{\large\ensuremath{\times}}}%
  {\textcolor{red!30!black}{#1}}\par
  % 正确句子行
  \noindent\makebox[1.8em][l]{\textcolor{blue!70!black}{\large\checkmark}}%
  {\textcolor{blue!30!black}{#2}}\par
  \vspace{0.3\baselineskip}% 紧凑间距（体现对照关系）
}

% 例句讲解命令
\newcommand{\explain}[3]{%
  \noindent \makebox[1em][l]{\raisebox{0.25ex}{\textcolor{lessonblue}{\large\ensuremath{\bullet}}}}%
  #1\par
  \noindent\hspace*{1.5em}{\small\textcolor{gray!30!black}{#2}}\par
  % ===== 语法讲解部分 (专业排版) =====
  \noindent\hspace*{1em}{%
    \color{lessonblue!80!black}%
    \small%
    % \raisebox{0.2ex}{\tiny\faLightbulbO}\ % 灯泡图标（可选）
    \color{lessonblue!60!black}%
    #3%
  }\par
  \vspace{0.4\baselineskip}%
}

% ==== 中文环境优化符号 ====
% ==== 音标 ====
% \newcommand{\phonetic}[1]{{\color{blue!70!black}/}{ #1}{{\color{blue!70!black}/}}}
\newfontfamily\ipafont{Charis SIL}[
  Path = fonts/,
  UprightFont = *-R,
  ItalicFont = *-I,
  BoldFont = *-B,
  BoldItalicFont = *-BI
]
\newcommand{\phonetic}[1]{{\color{blue!70!black}/\,}{\raisebox{0.25ex}{\ipafont\small\textcolor{blue!70!black}{#1}}}{{\color{blue!70!black}\,/}}}

% ==== 箭头 ====
\newcommand{\warr}{%
%   \hspace{0.15em}% 前间距
  \raisebox{-0.05ex}{\scalebox{0.9}{→}}% 提升0.2ex，缩小25%
  \hspace{0.25em}% 后间距
}

% ==== 缀字符 ====
\newcommand{\jie}[1]{%
  \raisebox{-0.1ex}{\scalebox{0.85}{–}}% 提升0.2ex，缩小15%
  \hspace{0.05em}% 微调间距
  #1%
}

% ===== 全局样式设置 =====
\ctexset{
  chapter={
    name={第,讲},
    number={\chinese{chapter}},
    format=\Huge\bfseries\color{lessonblue}\sffamily,
    aftername=\hspace{0.5em},
    beforeskip={2\baselineskip},
    afterskip={1.2\baselineskip}
  },
  section={
    format=\LARGE\bfseries\color{lessonblue!80!black}\sffamily,
    aftertitle={\hspace{0.3em}\color{lessonblue!50!black}·\hspace{0.3em}},
    beforeskip={1.2\baselineskip},
    afterskip={0.8\baselineskip}
  }
}